\clearpage
\phantomsection
\section*{Chương 6. Cửa sổ Hanning ROM $\times$ mẫu (Q1.15)}
\addcontentsline{toc}{section}{Chương 6. Cửa sổ Hanning ROM × mẫu (Q1.15)}
\pdfbookmark[1]{Chương 6. Cửa sổ Hanning ROM × mẫu}{chuong6}

\begin{learningobj}
Sau chương này bạn giải thích được vì sao khung âm thanh phải nhân cửa sổ Hanning trước khi tự tương quan; viết đúng công thức Praat ($N{+}1$) và đối chiếu dạng cổ điển ($N{-}1$); lượng tử hóa hệ số sang định dạng Q1.15; sinh tệp \texttt{hanning\_rom.hex} khớp bit với twin Python; hiện thực phép nhân ROM$\times$mẫu bằng một khối DSP trên Cyclone~V; và đọc được ngân sách M10K / DSP / ALM của khối cửa sổ.
\end{learningobj}

\subsection*{6.1. Hiện tượng thật: cắt khung không được ``cắt đứt'' sóng}

Giọng nói là một dòng sóng liên tục. Để đo chu kỳ rung dây thanh $T_0$, ta chỉ lấy một đoạn ngắn khoảng $40\,\text{ms}$ ($L = 640$ mẫu ở tần số lấy mẫu $f_s = 16\,\text{kHz}$), đúng như bản đồ khối ở Chương~0 đã ghi. Nếu lấy đoạn đó đúng kiểu cắt bằng kéo --- giữ nguyên biên độ đến mẫu cuối rồi nhảy đột ngột về 0 --- thì đoạn đó, nhìn từ miền tần số, giống một xung chữ nhật nhân với tín hiệu thật. Phép nhân trong miền thời gian tương đương \textbf{tích chập (convolution)} phổ tín hiệu với phổ của xung chữ nhật. Phổ xung chữ nhật là hàm $\mathrm{sinc}$: có một \textbf{thùy chính (mainlobe)} rộng và nhiều \textbf{thùy phụ (sidelobe)} cao khoảng $-13\,\text{dB}$ so với đỉnh (Chương~0 đã đổi $-13\,\text{dB}$ sang tỷ số biên độ $\approx 0{,}224$). Năng lượng ``rò'' qua các thùy phụ gọi là \textbf{rò phổ (spectral leakage)}.

Hệ quả với cao độ: đỉnh tự tương quan $\mathcal{R}_x[\tau]$ bị méo vì năng lượng giả từ các tần số láng giềng tràn sang đúng vùng tìm $F_0$ ($75 \ldots 600\,\text{Hz}$). Trực giác đời thường: bạn đang nghe một nốt đàn sạch, nhưng mỗi lần cắt đoạn lại ``vỗ'' hai đầu đoạn --- tiếng vỗ đó tạo thêm tiếng lách cách phổ, không có trong bản gốc.

Cửa sổ \textbf{Hanning} làm hai đầu khung về gần 0 một cách mượt, giống vuốt dần biên độ thay vì cắt đứt. Phổ Hanning có thùy phụ khoảng $-32\,\text{dB}$ (biên độ còn khoảng $1/40$ so với thùy chính, Chương~0). Thùy chính rộng hơn một chút (độ phân giải tần số kém hơn một ít), nhưng với bài đo $F_0$ bằng tự tương quan thời gian thì đánh đổi này rất đáng: ta cần đỉnh $\mathcal{R}_x$ sạch hơn là phổ sắc như dao.

\begin{physicalinsight}
Nhân cửa sổ $x[n]\cdot w[n]$ trên silicon chỉ là một phép nhân mỗi mẫu. Toàn bộ ``phép màu'' giảm rò phổ đã được nướng sẵn vào bảng hệ số $w[n]$. Không cần CORDIC, không cần sinh $\cos$ lúc chạy.
\end{physicalinsight}

\begin{figure}[H]
\centering
\begin{tikzpicture}[font=\small, >=Stealth]
  \draw[->] (0,0)--(8.2,0) node[right]{$n$};
  \draw[->] (0,0)--(0,2.8) node[above]{$w[n]$};
  \draw[thick, gray!70] (0.3,2.2)--(7.5,2.2) node[right,font=\scriptsize]{chữ nhật};
  \draw[thick, teal] (0.3,0)
    .. controls (1.5,0.15) and (2.2,2.2) .. (4.0,2.35)
    .. controls (5.8,2.2) and (6.5,0.15) .. (7.5,0);
  \node[teal, font=\scriptsize] at (4.0,2.6){Hanning};
  \draw[dashed, red!60] (0.3,0)--(0.3,2.2);
  \draw[dashed, red!60] (7.5,0)--(7.5,2.2);
  \node[red!70, font=\scriptsize, align=center] at (1.2,1.1){cắt đột\\gây rò phổ};
\end{tikzpicture}
\caption{Cửa sổ chữ nhật giữ biên độ đến sát biên rồi cắt đứt; Hanning vuốt hai đầu về 0, hạ thùy phụ từ khoảng $-13\,\text{dB}$ xuống khoảng $-32\,\text{dB}$.}
\label{fig:ch06_hanning_shape}
\end{figure}

\subsection*{6.2. Toán từng bước --- mọi ký hiệu trước khi dùng}

Gọi $L$ (hoặc $N$) là số mẫu trong một khung phân tích. Với sàn cao độ $75\,\text{Hz}$ ở $f_s = 16\,\text{kHz}$, Praat chọn khung dài khoảng ba chu kỳ: $L = 640$. Sau khi đã trừ mean cục bộ (Chương~3), tín hiệu cửa sổ là
\[
\tilde{x}[n] = (x[n] - \mu)\, w[n], \qquad n = 0,1,\ldots,L-1.
\]
Hai định nghĩa Hanning hay gặp --- phải chọn một và ghi vào twin:

\begin{enumerate}[leftmargin=1.4em]
\item \textbf{Praat ($N{+}1$), dùng trong golden \texttt{lab\_1\_3\_5.py}:}
\[
w[i] = 0{,}5 - 0{,}5\cos\Bigl(\frac{2\pi i}{N+1}\Bigr), \qquad i = 1,2,\ldots,N.
\]
Hai đầu lớn hơn 0 một chút (không đúng bằng 0). Với $N = 640$, mẫu số là $641$.
\item \textbf{Cổ điển / scipy ($N{-}1$):}
\[
w[n] = \frac12\Bigl(1 - \cos\frac{2\pi n}{N-1}\Bigr), \qquad n = 0,1,\ldots,N-1.
\]
Hai đầu đúng bằng 0; đối xứng $w[n] = w[N-1-n]$ sạch --- tiện lưu nửa bảng ROM.
\end{enumerate}

Gói FPGA này lấy dạng \textbf{Praat $N{+}1$} làm hợp đồng bit-exact với \texttt{hanning\_rom.hex} (đầy đủ 640 từ). Dạng cổ điển chỉ dùng trong ví dụ tính tay nhỏ và khi muốn minh họa tối ưu nửa bảng.

Lượng tử \textbf{Q1.15} (đã định nghĩa ở Chương~0: một bit dấu/phần nguyên, 15 bit phần thập phân; giá trị $1{,}0$ ánh xạ thành $2^{15} = 32768$, nhưng số có dấu 16-bit bão hòa tại $+32767$):
\[
w_q[n] = \mathrm{sat}_{16}\!\bigl(\mathrm{round}(w[n]\cdot 32768)\bigr) \in [0,32767]
\]
(với hệ số Hanning không âm). Mẫu vào $x_q$ cũng Q1.15. Tích rồi kéo về Q1.15 bằng làm tròn gần nhất:
\begin{align*}
y_{\mathrm{acc}} &= x_q \cdot w_q &&\text{(Q2.30, 32-bit có dấu)},\\
y_q &= \mathrm{sat}_{16}\!\bigl((y_{\mathrm{acc}} + 2^{14}) \ggg 15\bigr).
\end{align*}
Cộng $2^{14} = 16384$ trước khi dịch phải 15 bit chính là round-to-nearest trên số học hai's complement khi dùng arithmetic shift.

Năng lượng cửa sổ (dùng ở Chương~8 để chuẩn hóa):
\[
\mathcal{R}_w[0] = \sum_{n=0}^{N-1} w_q[n]^2
\]
tính offline trên int64; twin in giá trị này vào header hex để đối chiếu.

\subsection*{6.3. Ví dụ tính tay --- mọi số đã chạy Python}

Mọi con số dưới đây được kiểm chứng bởi \texttt{verify/verify\_chuong\_06\_09.py}.

\paragraph{A. Dạng cổ điển $N=8$ (đối xứng nửa bảng).}
$w[n] = \tfrac12\bigl(1 - \cos(2\pi n/7)\bigr)$, $n = 0..7$, lượng tử bằng $\mathrm{round}(w\cdot 32767)$ để minh họa hệ số dương:
\begin{center}\small
\begin{tabular}{@{}rrl@{}}\toprule
$n$ & $w[n]$ float & $w_q$ \\
\midrule
0 & 0.000000 & 0 \\
1 & 0.188255 & 6169 \\
2 & 0.611260 & 20029 \\
3 & 0.950484 & 31145 \\
\bottomrule
\end{tabular}
\end{center}
Đối xứng: $w_q[7] = w_q[0]$, $w_q[6] = w_q[1]$, \ldots --- ROM vật lý chỉ cần 4 từ. Nhân: $x_q = 10000$, $w_q = 20029$ cho tích $200\,290\,000$; làm tròn $({\rm acc}+2^{14})\ggg 15 = 6112$. Đối chiếu float $10000\times 0{,}611260 = 6112{,}6$ --- khớp trong 1 LSB.

\paragraph{B. Dạng Praat $N=8$ (khớp tinh thần golden).}
$w[i] = 0{,}5 - 0{,}5\cos(2\pi i/9)$, $i = 1..8$, lượng tử $\mathrm{round}(w\cdot 32768)$ rồi bão hòa:
\begin{center}\small
\begin{tabular}{@{}rrl@{}}\toprule
$i$ & $w$ float & $w_q$ \\
\midrule
1 & 0.116978 & 3833 \\
2 & 0.413176 & 13539 \\
3 & 0.750000 & 24576 \\
4 & 0.969846 & 31780 \\
\bottomrule
\end{tabular}
\end{center}
Hai đầu không về 0 tuyệt đối --- đúng hành vi Praat.

\paragraph{C. Khung thật $N=640$ Praat.}
Các mốc Q1.15: $w_q[0] = 1$, $w_q[1] = 3$, $w_q[319] = w_q[320] = 32767$, $w_q[638] = 3$, $w_q[639] = 1$. Tổng bình phương $\mathcal{R}_w[0] = \sum w_q^2 = 258\,100\,276\,868$ (int64).

\subsection*{6.4. Ý nghĩa phần cứng --- ROM On-Chip Memory + DSP}

Cửa sổ Hanning \emph{không} cần FIR Compiler~II: hệ số bất biến với $L$ cố định. Chọn IP \textbf{On-Chip Memory} kiểu ROM (hoặc RAM 1-PORT khởi tạo \texttt{.hex}/\texttt{.mif}):
\begin{itemize}[leftmargin=1.4em]
\item Width $= 16$ (Q1.15), depth $= 640$ (đầy đủ, khớp Praat) hoặc $320$ nếu dùng dạng cổ điển nửa bảng.
\item Registered output: 1 chu kỳ đọc --- FSM phải chờ đúng latency này.
\item Một khối \textbf{DSP} (\emph{Digital Signal Processing} block) trên Cyclone~V nhân $18\times 19$ mỗi mẫu: $x_q\times w_q$.
\end{itemize}
Không gọi FIR Compiler chỉ để nhân cửa sổ: lãng phí logic cấu hình systolic. Đổi $L$ hoặc công thức $w$ thì \textbf{bắt buộc rebuild hex} và nạp lại ROM; quên bước này thì $F_0$ sai im lặng.

\paragraph{RTL datapath --- giải thích từng dòng.}
\begin{verbatim}
module window_hanning #(
  parameter N = 640
)(
  input  wire               clk, rst_n,
  input  wire               in_valid,
  output wire               in_ready,
  input  wire signed [15:0] x_q,        // Q1.15
  output reg                out_valid,
  input  wire               out_ready,
  output reg  signed [15:0] y_q         // Q1.15 windowed
);
  (* ram_init_file = "hanning_rom.hex" *)
  reg [15:0] rom [0:N-1];
  reg [9:0]  n;                         // 0..N-1
  wire signed [15:0] w_q = rom[n];
  wire signed [31:0] prod = x_q * w_q;  // Q2.30
  wire signed [31:0] rounded = prod + 32'sd16384;
  wire signed [31:0] shifted = rounded >>> 15;
  function automatic signed [15:0] sat16;
    input signed [31:0] v;
    begin
      if (v > 32'sd32767) sat16 = 16'sd32767;
      else if (v < -32'sd32768) sat16 = -16'sd32768;
      else sat16 = v[15:0];
    end
  endfunction
  assign in_ready = out_ready | ~out_valid;
  always @(posedge clk or negedge rst_n) begin
    if (!rst_n) begin n <= 0; out_valid <= 0; y_q <= 0; end
    else if (in_valid && in_ready) begin
      y_q <= sat16(shifted);
      out_valid <= 1'b1;
      n <= (n == N-1) ? 10'd0 : n + 10'd1;
    end else if (out_valid && out_ready) out_valid <= 1'b0;
  end
endmodule
\end{verbatim}
\texttt{in\_ready}/\texttt{out\_ready} là handshake Avalon-ST (sink/source, Chương~0): khi downstream chưa sẵn sàng, upstream phải giữ \texttt{valid}. Hàm \texttt{sat16} bão hòa thay vì wrap, tránh đảo dấu khi tích lớn.

Nếu muốn tiết kiệm một nửa BRAM với dạng cổ điển: địa chỉ \texttt{(n < 320) ? n : (639-n)} và ROM depth 320. Twin phải cùng công thức --- không trộn Praat hex với địa chỉ nửa bảng cổ điển.

\subsection*{6.5. Twin Python bit-exact}

\begin{verbatim}
import numpy as np
Q15 = 1 << 15

def hanning_praat_q15(N=640):
    i = np.arange(1, N + 1, dtype=np.float64)
    w = 0.5 - 0.5 * np.cos(2.0 * np.pi * i / (N + 1))
    return np.clip(np.round(w * Q15), 0, Q15 - 1).astype(np.int64)

def apply_window(xq, wq):
    prod = xq.astype(np.int64) * wq.astype(np.int64)
    return ((prod + (1 << 14)) >> 15).astype(np.int64)

wq = hanning_praat_q15(640)
with open("hex/hanning_rom.hex", "w") as f:
    for v in wq:
        f.write(f"{int(v) & 0xFFFF:04X}\n")
print("Rw0 =", int(np.sum(wq * wq)))
\end{verbatim}
Mọi thay đổi round/sat phải sửa \textbf{cùng chỗ} RTL và twin --- đó là hợp đồng bit-exact. Script kiểm số tay: \texttt{verify/verify\_chuong\_06\_09.py}.

\subsection*{6.6. Chi phí silicon --- viết tắt mở rộng lần đầu trong chương}

\begin{costbox}
\textbf{Ngân sách khối cửa sổ.}
\begin{itemize}[leftmargin=1.2em]
\item \textbf{M10K} (khối BRAM nhúng $10\,240$ bit trên Cyclone~V, Chương~0): ROM đầy đủ $640\times 16 = 10\,240$\,bit $=$ đúng 1 M10K. Nửa bảng cổ điển $320\times 16 = 5\,120$\,bit vẫn chiếm 1 M10K (không chia nhỏ hơn một block).
\item \textbf{DSP}: 1 nhân $18\times 19$ mỗi mẫu cửa sổ (có thể dùng chung với MAC ở Pha~2 nếu FSM tuần tự).
\item \textbf{ALM} (\emph{Adaptive Logic Module} --- đơn vị logic cơ bản Cyclone~V, gồm LUT look-up table + flip-flop): con trỏ $n$ + sat + handshake $\lesssim 80$ ALM.
\item \textbf{Fmax} (tần số đồng hồ tối đa sau place-and-route): $\gg 50\,\text{MHz}$ --- đường nhân một chu kỳ không phải nút thắt.
\item \textbf{Latency}: 1 chu kỳ đọc ROM (cộng 0/1 chu kỳ nhân pipeline tùy suy diễn).
\item Vì sao ROM chứ không CORDIC: $w[n]$ bất biến; tính offline một lần rẻ hơn sinh $\cos$ lúc chạy hàng nghìn lần mỗi giây.
\end{itemize}
\end{costbox}

\subsection*{6.7. Bài tập và lời giải giải thích}

\paragraph{Bài tập 6.1.} Sinh hex Praat $N=8$; đối chiếu bốn từ đầu với bảng mục 6.3.B.

\paragraph{Bài tập 6.2.} RTL sim: khung toàn mẫu $32767$ nhân Hanning --- mô tả biên độ $y$ ở hai đầu và giữa khung. Giải thích vì sao.

\paragraph{Bài tập 6.3.} Cố ý dùng hex Praat nhưng địa chỉ nửa bảng cổ điển: $\mathcal{R}_x[\tau]$ lệch twin thế nào?

\paragraph{Bài tập 6.4.} Tính tay $x_q = -10000$, $w_q = 20029$ với cùng quy tắc $+16384$ rồi $\ggg 15$. Kết quả?

\noindent\textbf{Lời giải giải thích.}
\begin{enumerate}[leftmargin=1.4em]
\item Kỳ vọng bốn từ đầu $3833$, $13539$, $24576$, $31780$ (script đã assert).
\item Giữa khung $w\approx 1$ nên $y\approx 32767$; hai đầu $w$ nhỏ nên $y$ nhỏ --- đúng hình quả núi Hanning. Đây là hệ quả trực tiếp của công thức Praat, không phải lỗi bão hòa.
\item Hệ số sai đối xứng làm méo $\mathcal{R}_x$, đỉnh $\tau$ lệch vài mẫu, dẫn đến sai vài chục cent khi đối chiếu Praat ở Chương~13.
\item Tích $-200\,290\,000$; cộng $16384$ rồi arithmetic $\ggg 15$ cho $-6112$ (đối xứng với ví dụ dương ở mục 6.3.A). Script xác nhận \texttt{q15\_mul(-10000, 20029) == -6112}.
\end{enumerate}
